\section{Problem Statement}
\label{sec:problem-statement}

The preceding section motivates a gap between available observations and the
model needed to profile logical task attempts in queue-worker systems. This
section states the problem addressed by the paper. It defines the system model,
available information, expected outputs, main constraints, assumptions, and
non-goals. The purpose is to state the problem without introducing the full model
that later sections define.

At a high level, the problem is to estimate the behavior of a task attempt from
its task identity, intrinsic task features, execution-context information, and
available history while respecting profiling overhead. The estimate must remain
structured. It must preserve active resource demand, worker occupancy,
uncertainty, and operational risk as distinct outputs because these quantities
support different decisions.

\subsection{System Model}
\label{sec:problem-system-model}

We consider queue-worker systems in which producers submit task instances to one
or more queues and workers execute task attempts. A task kind identifies the
semantic type of work requested. A task instance is one queued unit of that kind.
A task attempt is one concrete execution of a task instance. A single task
instance may produce multiple attempts because of retries, lease expiration,
timeouts, cancellation, worker failure, or rescheduling.

The profiling boundary is the task attempt. This boundary is narrower than a
worker process, container, node, or worker pool, but broader than a single
function call or trace span. An attempt may include sequential, parallel,
conditional, retried, or partially overlapping phases. The paper does not assume
that every phase is perfectly instrumented.

Each attempt executes under an execution context. Execution context is the
input-side state surrounding the attempt; it is not the cost of the attempt. It
may include worker state, worker-pool pressure, node pressure, dependency
health, network condition, disk pressure, cache state, replica state, tenant
pressure, retry state, and domain-specific hidden state. Some context factors
are directly observable, some are derived from raw observations, and some remain
hidden.

The model is task-kind-specific. Different task kinds may require different
features, observation points, and profiling policies. A replication task may be
sensitive to replica health and network condition. A compute-oriented transform
may be more sensitive to record count, algorithm mode, CPU pressure, allocation
behavior, and memory pressure. The problem therefore cannot assume one global
feature set or one global metric list that is equally relevant to all task
kinds.

\subsection{Available Information}
\label{sec:problem-available-information}

For an attempt $a$, the profiler may have access to four classes of information.

First, task identity describes what is being executed. It may include the task
kind, task version, queue, producer, tenant, priority, deadline, and attempt
number. Task identity selects the relevant task-kind-specific model; it is not
by itself a cost estimate.

Second, intrinsic task features describe work contained in the task instance.
Examples include payload size, record count, object count, partition size, batch
size, input format, operation mode, declared resource hints, or domain metadata.
These features may be available before execution and can provide a baseline for
expected work.

Third, execution-context information describes the environment in which the
attempt is executed or expected to execute. Examples include worker-pool
pressure, node CPU or memory pressure, disk queue depth, network instability,
dependency health, connection-pool pressure, cache warmth, replica health,
replication lag, tenant quota usage, and retry state. These values are
input-side conditions. They may influence resource demand or worker occupancy,
but they are not themselves worker-occupancy values.

Fourth, historical observations describe previous attempts of the same or
related task kinds. They may include observed resource usage, observed worker
occupancy, durations, retry counts, outcomes, residuals, tail quantiles, or drift
indicators. Historical data may be unavailable for new task kinds, new versions,
rare contexts, or recently changed workloads. The problem must therefore support
both calibrated and cold-start cases.

Raw observations are not assumed to be model-ready. Metrics, logs, traces,
runtime profiles, and domain events must be interpreted, normalized, and mapped
into canonical features or observed output values before the model can use them.

\subsection{Desired Outputs}
\label{sec:problem-desired-outputs}

The profiler should produce a structured estimate for a task attempt or for a
class of comparable attempts.

The first output is active resource demand. For each relevant resource dimension
$r \in \mathcal{R}$, the model should estimate $\widehat{D}_r(a)$ or a
distribution over possible values. Relevant dimensions may include CPU time,
memory peak, disk reads, disk writes, network input, network output, allocation
volume, or accelerator usage.

The second output is worker occupancy. For each relevant occupancy dimension
$o \in \mathcal{O}$, the model should estimate $\widehat{U}_o(a)$ or a
distribution over possible values. Relevant dimensions may include worker-slot
time, lease time, connection hold time, lock hold time, semaphore hold time, and
tenant concurrency time. Occupancy is modeled separately because an attempt can
hold scarce execution capacity without consuming much active CPU or memory.

The third output is uncertainty and confidence. The model should indicate how
reliable its estimates are, whether residual error is expected to be high,
whether behavior has drifted, whether historical samples are insufficient, and
whether tail behavior is operationally significant. The paper does not require a
single universal uncertainty metric.

The fourth output is operational risk. Risk is not the same as task cost. It
captures the consequence of underestimating, misclassifying, or mishandling an
attempt. Relevant risks may include queue starvation, retry amplification, SLO
violation, replication lag, tenant impact, dependency overload, and cascading
overload.

The fifth output is a profiling-policy recommendation. Given the task kind,
context, uncertainty, risk, and profiling budget, the system should determine
what observation strategy is justified. The policy output is a profiling
decision, not a final scheduling, routing, or placement decision.

\subsection{Core Challenges}
\label{sec:problem-core-challenges}

The first challenge is conditional behavior. Attempts with the same task kind
and similar intrinsic features may behave differently under different execution
contexts. A model based only on task kind, queue, or payload bucket can provide a
baseline, but it cannot explain context-dependent amplification, suppression, or
drift.

The second challenge is multidimensionality. Active resource demand, bounded
worker occupancy, uncertainty, and operational risk have different meanings and
units. Combining them prematurely into one scalar hides information needed by
downstream control layers. A scalar score may be useful, but it should be a
projection of a structured estimate rather than a replacement for it.

The third challenge is attribution. Process-level, container-level, node-level,
and queue-level observations are often easier to collect than
task-attempt-level observations. Aggregate observations may indicate pressure,
but they do not automatically identify which task kind, instance, attempt,
phase, context factor, or dependency produced the observed behavior.

The fourth challenge is semantic interpretation. Raw observations do not state
their model role. Without explicit semantic roles, a model may double-count
related signals or treat symptoms as independent causes. The problem therefore
requires a mapping from raw observations to canonical features and observed
output values.

The fifth challenge is compositional behavior. A task kind may include CPU
processing, network transfer, dependency calls, retries, commits, and cleanup.
Some quantities compose by summation, others by maximum, critical path, retry
expansion, or occupancy-specific rules. The profiling problem must preserve
enough structure to combine component behavior consistently.

The sixth challenge is partial observability. Important influences may be
missing, stale, hidden, or domain-specific. The model cannot assume that all
relevant state is observable. Unexplained behavior should remain visible through
residual error, reduced confidence, or a higher profiling level.

The seventh challenge is profiling overhead. Instrumentation is not free.
Detailed tracing, runtime profiling, high-cardinality labels, frequent sampling,
and telemetry export can affect the workload being measured. A profiling model
must decide not only what would be useful to observe, but what is worth observing
under an overhead budget.

\subsection{Scope and Assumptions}
\label{sec:problem-scope-assumptions}

This paper assumes that task attempts can be identified at least by task kind
and attempt boundary. The model is less useful when all work is invisible inside
one undifferentiated worker loop. The paper does not require perfect
instrumentation, but it assumes that minimal accounting is possible for at least
some attempts.

The paper assumes that some task metadata or intrinsic features may be available
before or during execution. The exact features are domain-specific. A validation
task, replication task, external API fan-out task, and media transformation task
need not expose the same feature set.

The paper assumes that execution-context information may be available but
incomplete. It may come from workers, nodes, queues, dependencies, telemetry
systems, domain services, or historical profiles. The model must tolerate
missing, stale, noisy, or partially aggregated context.

The paper assumes that profiling decisions are useful even when final
scheduling, routing, placement, or admission decisions are made by another layer.
The profiling model supplies structured estimates and observation policies.
Downstream control layers may consume these outputs, but their full design is
outside the scope of the paper.

The paper also assumes an open-world metric environment. No core system can know
all possible domain-specific metrics in advance. The model must therefore allow
raw observations to be mapped into canonical features and must keep missing or
unexplained influences explicit instead of hiding them behind a fixed global
metric list.

\subsection{Non-Goals}
\label{sec:problem-non-goals}

This paper does not propose a complete distributed scheduler. It does not define
a placement algorithm, routing algorithm, gossip protocol, load-balancing policy,
or cluster-wide resource allocator. Such systems may use the estimates produced
by task profiling, but they are not the object of this paper.

This paper does not replace metrics, logs, traces, runtime profiles, or existing
observability tools. These systems provide observations. The problem addressed
here is how to interpret observations for task-kind-specific profiling and how
to decide which observations are justified.

This paper does not claim that Task Behavior Graphs prove causality. A graph can
make modeled influence assumptions explicit, but causal validation requires
additional evidence, experiments, interventions, or domain knowledge. The graph
is a profiling model, not a causal proof.

This paper does not attempt to define a globally complete taxonomy of all
workload metrics. It assumes that new domains can introduce new observations,
features, and behavior components. Completeness is task-kind- and
domain-specific rather than global.

This paper does not require exact prediction for arbitrary tasks. The goal is to
produce structured estimates, expose uncertainty, and guide profiling policy.
High residual error, drift, missing features, or high tail risk are valid model
outputs, not failures to hide.

Finally, this paper does not specify a full production implementation of a
package manager, expression language, runtime virtual machine, or cross-language
manifest compiler. Later sections discuss manifests and implementation
considerations only to the extent needed to make the proposed model concrete.
